\chapter{How to Work Through the Book}

This appendix is a practical guide for self-study.  The manuscript is long because the subject has many layers, but the working rhythm is simple.  For each section, copy the central derivation by hand, mark the boundary condition or normalization being used, and then solve one exercise without looking at the answer.  The goal is not speed.  The goal is to make each symbol accountable.

\section{A derivation log}

Keep a derivation log with four columns: formula, assumption, boundary condition, and check.  The formula is the line of mathematics being used.  The assumption records what is held fixed or treated as small.  The boundary condition records why a surface term vanishes or why it remains.  The check is a simple limit, dimensional test, or special case.

\section{When to move on}

Move on from a section when you can reproduce its central derivation without the prose.  You do not need to remember every symbol, but you should know why each derivative, commutator, source, or integral appears.  If a derivation still feels like a row of decorations, return to the finite model at the beginning of the section.
